\chapter{Il pianeta che bolle}

\begin{flushright}
	\textit{``Il riscaldamento globale \`{e} il perfetto nemico del cervello umano: troppo lento, troppo complesso, troppo astratto.''}
\end{flushright}

\bigskip
\bigskip

Se volessimo progettare un problema su misura per mettere in scacco tutti i bias cognitivi umani contemporaneamente, probabilmente inventeremmo qualcosa di molto simile al cambiamento climatico. \`{E} come se fosse stato concepito apposta per sfruttare ogni debolezza del nostro cervello antico. E infatti, nel 2024, la differenza di opinioni sul clima tra elettori repubblicani e democratici americani ha toccato il record storico di 38 punti percentuali. Non stiamo pi\`{u} parlando di scienza. Stiamo assistendo al trionfo dei nostri bias ancestrali su scala planetaria.

\section{Il bias del presente contro il futuro remoto}

Il nostro cervello \`{e} ottimizzato per minacce immediate e concrete: il predatore dietro il cespuglio, il rivale armato di clava. Il cambiamento climatico \`{e} l'opposto di tutto questo: lento, graduale, statistico, proiettato in un futuro che sembra sempre troppo lontano per meritare attenzione oggi. \`{E} come chiedere a chi ha imparato a guidare su strade di montagna di pilotare un transatlantico: tecnicamente sono entrambi veicoli, ma richiedono competenze e tempi di reazione completamente diversi.

I dati confermano tutto questo in modo spietato: l'88\% dei democratici americani crede che il riscaldamento globale sia in corso, contro il 54\% dei repubblicani. Ma il dato pi\`{u} rivelatore \`{e} un altro: anche tra persone che vivono nelle stesse regioni, colpite dagli stessi eventi climatici estremi, la percezione diverge radicalmente in base all'appartenenza politica. L'86\% dei democratici afferma di vedere il cambiamento climatico colpire la propria comunit\`{a}; solo il 41\% dei repubblicani riconosce gli stessi impatti, nella stessa citt\`{a}, sotto lo stesso cielo sempre pi\`{u} caldo. Non \`{e} cecit\`{a} fisica: \`{e} il cervello che filtra la realt\`{a} attraverso le lenti del proprio gruppo di appartenenza.

Su questo scarto tra ci\`{o} che il cervello sa fare bene e ci\`{o} che il clima gli richiede, un filosofo tedesco ha scritto forse la riflessione pi\`{u} lucida del secolo scorso. Hans Jonas, nel suo \textit{Principio responsabilit\`{a}} del 1979, part\`{i} da una constatazione semplice: l'etica tradizionale, da Aristotele in poi, era stata pensata per azioni dagli effetti immediati e circoscritti, tra persone che si conoscono e che vivono nello stesso presente.\footnote{Jonas, H. (1979). \textit{Das Prinzip Verantwortung}. Insel Verlag, Francoforte.} Ma la tecnologia moderna, sosteneva Jonas, ha dato all'uomo un potere del tutto nuovo: quello di produrre conseguenze che si estendono su scala planetaria e temporale, capaci di danneggiare persone che non sono ancora nate e che non potranno mai avere voce in capitolo sulle nostre decisioni di oggi. Di fronte a un potere cos\`{i} sproporzionato rispetto alla nostra capacit\`{a} di prevederne gli effetti, Jonas propose quella che chiam\`{o} l'\index[main]{euristica della paura}euristica della paura: dobbiamo lasciare che l'immagine dei possibili disastri futuri, anche solo temuti e non certi, orienti le nostre azioni presenti, invece di aspettare la prova definitiva prima di agire. Non serve la certezza assoluta che il peggio accadr\`{a}: basta la possibilit\`{a} che accada, quando la posta in gioco \`{e} l'esistenza stessa delle generazioni future. \`{E} l'esatto contrario di come il nostro cervello \`{e} biologicamente predisposto a funzionare, ottimizzato com'\`{e} per reagire a ci\`{o} che \`{e} certo e immediato, non a ci\`{o} che \`{e} probabile e remoto. La sua euristica della paura \`{e}, in un certo senso, un correttivo filosofico che dovrebbe compensare esattamente il bias con cui abbiamo aperto questa sezione.

\section{L'ottimismo irrazionale}

Tendiamo sistematicamente a sottostimare i rischi futuri e a sovrastimare la nostra capacit\`{a} di affrontarli quando si presenteranno. \`{E} l'ottimismo bias applicato su scala planetaria: ce la caveremo, la tecnologia ci salver\`{a}. Questo ottimismo \`{e} stato evolutivamente vantaggioso, perch\'{e} gli antenati troppo pessimisti tendevano a paralizzarsi, ma diventa un problema quando si tratta di minacce che richiedono un'azione preventiva su larga scala. Il 76\% degli americani crede che la tecnologia avanzata offra speranza per il futuro, e il 55\% \`{e} convinto che risolver\`{a} la maggior parte dei problemi climatici entro 50 anni. C'\`{e} per\`{o} un paradosso crudele in questo ottimismo tecnologico: pi\`{u} crediamo che la tecnologia ci salver\`{a} domani, meno agiamo oggi, esattamente come si rimanda una dieta perch\'{e} tanto prima o poi inventeranno una pillola dimagrante.

Questa fiducia illimitata nella tecnologia futura, unita all'incapacit\`{a} di temere davvero le conseguenze della tecnologia presente, \`{e} precisamente ci\`{o} che il filosofo Günther Anders analizz\`{o} gi\`{a} negli anni Cinquanta, in un'opera dal titolo quasi profetico, \textit{L'uomo \`{e} antiquato}.\footnote{Anders, G. (1956). \textit{Die Antiquiertheit des Menschen}. C. H. Beck, Monaco.} Anders sosteneva che l'umanit\`{a} soffre di una \index[main]{cecit\`{a} apocalittica}cecit\`{a} apocalittica strutturale: siamo psicologicamente incapaci di immaginare per intero le catastrofi che le nostre stesse tecnologie rendono possibili, perch\'{e} la nostra immaginazione morale resta tarata su una scala umana, non su quella planetaria dei sistemi che abbiamo costruito. Ma il concetto pi\`{u} sorprendente che introdusse fu quello della \textit{prometheische Scham}, la vergogna prometeica: il senso di inferiorit\`{a} che l'uomo prova di fronte alle proprie macchine, cos\`{i} perfette ed efficienti rispetto alla fragilit\`{a} biologica di chi le ha create. \`{E} una vergogna che genera una fiducia quasi religiosa nella tecnologia, come se solo essa, e non pi\`{u} l'azione umana diretta, potesse davvero risolvere i problemi che l'azione umana ha creato. L'ottimismo tecnologico descritto sopra, il 55\% convinto che la tecnologia risolver\`{a} tutto entro qualche decennio, \`{e} quasi un caso da manuale di questa vergogna prometeica: deleghiamo alla macchina la fiducia che non riusciamo pi\`{u} a riporre in noi stessi, e nel farlo ci esentiamo dall'agire nel presente.

Il caso americano lo dimostra con chiarezza. L'amministrazione Biden aveva investito 370 miliardi di dollari nell'\textit{Inflation Reduction Act}, il pi\`{u} grande investimento climatico della storia. Nel gennaio 2025, primo giorno della nuova amministrazione, gran parte di quelle misure fu cancellata con un ordine esecutivo dal nome esplicito, ``liberare l'energia americana'': un ritorno, in sostanza, a bruciare tutto ci\`{o} che si pu\`{o} trovare.

\section{Il bias dell'identit\`{a} tribale}

Il cambiamento climatico \`{e} diventato un marcatore di identit\`{a} politica, e questo attiva tutti i nostri bias di gruppo. Non valutiamo pi\`{u} le evidenze scientifiche in base alla loro validit\`{a}, ma in base a quale gruppo le sostiene: se la mia parte politica non crede al riscaldamento globale, tender\`{o} a non crederci nemmeno io, indipendentemente dalle prove. Nel biennio 2024 e 2025 questo tribalismo ha raggiunto livelli quasi paradossali: la fiducia negli scienziati del clima varia dall'88\% tra i democratici al 43\% tra i repubblicani, un divario di 45 punti percentuali. Solo il 12\% dei repubblicani considera il clima una priorit\`{a}, collocandolo ultimo tra venti questioni politiche; per i democratici la percentuale sale al 59\%. I ricercatori chiamano questo fenomeno \index[cognitive]{cognizione identity protective}cognizione identity protective: quando l'appartenenza politica viene resa saliente nella mente di una persona, la sua fiducia nel cambiamento climatico si riduce automaticamente, quasi indipendentemente dai fatti che ha davanti.

\section{L'algoritmo dell'apocalisse}

C'\`{e} per\`{o} un fattore che i nostri antenati non avrebbero mai potuto immaginare: l'ecosistema digitale della disinformazione, capace di amplificare questi bias su scala industriale. Un esperimento ha simulato l'esperienza di un utente scettico sul clima su un social network, e l'algoritmo ha raccomandato pagine contenenti disinformazione climatica pura nel giro di pochi clic.\footnote{Global Witness. Report sperimentale sulla disinformazione climatica algoritmica.} Otto delle dieci trasmissioni online pi\`{u} popolari diffondono sistematicamente contenuti fuorvianti sul clima, e alcune di queste sono finanziate con milioni di dollari da interessi legati ai combustibili fossili. Durante le conferenze internazionali sul clima del 2023 e del 2024 sono stati identificati centinaia, in un caso quasi duemila, account falsi che promuovevano narrazioni favorevoli ai combustibili fossili. E la maggioranza della disinformazione climatica pi\`{u} recente ha abbandonato la negazione diretta per tattiche pi\`{u} sofisticate: non pi\`{u} ``non sta succedendo'', ma ``le soluzioni non funzionano'' o ``il cambiamento climatico ha anche dei benefici''. \`{E} l'evoluzione naturale del negazionismo, che si adatta pur di continuare a offrire al cervello ci\`{o} che desidera sentirsi dire.

Dietro questa macchina c'\`{e} una logica economica precisa. Le principali compagnie petrolifere hanno speso centinaia di milioni di dollari in attivit\`{a} di lobbying negli ultimi quindici anni. Alcune di esse conoscevano gi\`{a} dagli anni Settanta, attraverso i propri modelli interni, i rischi climatici che oggi vediamo realizzarsi, ma scelsero di investire una frazione minima delle proprie spese in tecnologie a basse emissioni, dedicando il resto a campagne per confondere l'opinione pubblica. \`{E} il bias dello status quo, l'inerzia verso ci\`{o} che gi\`{a} conosciamo e su cui abbiamo gi\`{a} investito, moltiplicato per miliardi di dollari e trasformato in strategia industriale.

\section{La frattura generazionale}

C'\`{e} per\`{o} una crepa in questo muro. Il 40\% dei repubblicani sotto i 45 anni crede ormai che il cambiamento climatico sia causato principalmente dall'uomo, contro il 26\% del 2017. La generazione pi\`{u} giovane \`{e} anche quella pi\`{u} colpita nei fatti: il 61\% ha gi\`{a} vissuto direttamente alluvioni, siccit\`{a} o problemi di accesso all'acqua potabile negli ultimi due anni, e due terzi si aspettano che la situazione peggiori nel corso della propria vita. \`{E}, letteralmente, la prima generazione della storia che sa per certo di ereditare un pianeta in condizioni peggiori di quello dei propri genitori.

Questa frattura tra generazioni tocca una domanda antica quanto la filosofia politica: cosa dobbiamo, esattamente, a chi verr\`{a} dopo di noi? Nel 1790, in polemica con gli eccessi della Rivoluzione francese, il pensatore irlandese Edmund Burke scrisse una delle definizioni pi\`{u} durature della societ\`{a} umana nelle sue \textit{Riflessioni sulla rivoluzione in Francia}: la societ\`{a}, sosteneva, \`{e} un contratto, ma non un contratto qualsiasi, stipulato solo tra chi \`{e} vivo in un dato momento.\footnote{Burke, E. (1790). \textit{Reflections on the Revolution in France}. J. Dodsley, Londra.} \`{E}, scriveva, un patto tra chi vive, chi \`{e} gi\`{a} morto e chi deve ancora nascere: ogni generazione riceve in eredit\`{a} un mondo che non ha costruito da sola, ed \`{e} tenuta a consegnarlo, migliorato o almeno non deteriorato, a quella successiva. Burke lo scriveva pensando alle istituzioni politiche e alla continuit\`{a} sociale, non al clima, ma la sua intuizione si applica qui con una precisione quasi impressionante: la frattura generazionale che osserviamo oggi non \`{e} soltanto un divario di opinioni tra genitori e figli, \`{e} la rottura concreta di quel patto tra generazioni, in cui chi ha ricevuto un pianeta stabile lo sta restituendo instabile, senza aver mai chiesto il permesso a chi verr\`{a} dopo. E il paradosso \`{e} amaro: chi meno ha contribuito al problema, i giovani, i gruppi a basso reddito, le comunit\`{a} pi\`{u} esposte, \`{e} anche chi ha meno potere per correggerne la rotta.

\section{Il paradosso della rana bollita}

C'\`{e} una vecchia immagine, scientificamente falsa ma metaforicamente efficace, di una rana immersa in un'acqua che si scalda cos\`{i} gradualmente da non farla mai saltare fuori, finch\'{e} non bolle senza accorgersene. Il cambiamento climatico \`{e} la pentola, noi siamo la rana, e i nostri bias sono il coperchio che ci impedisce di reagire in tempo. Il cervello non \`{e} equipaggiato per percepire minacce che si manifestano nell'arco di decenni invece che di secondi: un aumento di un grado e mezzo in cinquant'anni non genera alcun allarme biologico, mentre un ruggito nella savana produceva adrenalina immediata. Ecco perch\'{e} gli eventi climatici estremi, uragani, alluvioni, ondate di calore, sono paradossalmente utili: trasformano l'astratto in concreto, il futuro in presente. Ma anche qui interviene il bias tribale: se il proprio gruppo ha gi\`{a} deciso che quegli eventi sono ``normale variabilit\`{a} meteorologica'', si continuer\`{a} a vedere soltanto tempo brutto, mai un collasso in corso.

\`{E} precisamente questa la cecit\`{a} apocalittica descritta da Anders: non l'ignoranza dei fatti, che spesso conosciamo perfettamente, ma l'incapacit\`{a} strutturale di lasciare che quei fatti producano nella mente un allarme proporzionato alla loro gravit\`{a}. Sappiamo che l'acqua si sta scaldando. Ma il nostro sistema di allarme, tarato su pericoli immediati, continua a sussurrarci che va tutto bene.

\section{Una soluzione parziale}

Come si esce, allora, da questa trappola cognitiva quasi perfetta? La verit\`{a} scomoda \`{e} che probabilmente non se ne esce del tutto: i bias che ci impediscono di affrontare il cambiamento climatico sono gli stessi che ci hanno tenuto vivi per milioni di anni, e non si possono disattivare a comando. Esistono per\`{o} strategie parziali. I dati mostrano che il sostegno per certe politiche climatiche resta trasversale agli schieramenti quando vengono inquadrate nel modo giusto: gli sgravi fiscali per l'efficienza energetica raccolgono il 75\% di consenso tra gli elettori repubblicani, e le tecnologie di cattura del carbonio, se presentate come innovazione di mercato anzich\'{e} come imposizione ambientale, arrivano al 69\% di approvazione conservatrice. \`{E} come parlare al cervello antico nella sua stessa lingua: non ``salviamo il pianeta'', un concetto troppo astratto per attivare alcunch\'{e}, ma ``rendiamo la nazione indipendente dal punto di vista energetico'', un obiettivo insieme tribale e concreto. Anche la resilienza a livello locale offre una via percorribile: decine di stati americani mantengono impegni climatici autonomi indipendentemente dalle oscillazioni della politica federale, applicando in scala ridotta il principio jonassiano secondo cui, di fronte a un rischio collettivo enorme, conviene agire anche senza attendere un consenso totale, perch\'{e} l'inazione stessa \`{e} gi\`{a} una scelta, e generalmente la peggiore.

La grande ironia di questa storia \`{e} che i bias che ci impediscono di affrontare il cambiamento climatico sono perfettamente razionali dal punto di vista evolutivo: il cervello che ignora le minacce lente e astratte per concentrarsi su benefici immediati e concreti ha dominato per millenni. Il problema non \`{e} il cervello in s\'{e}. Il problema \`{e} che abbiamo costruito una civilt\`{a} capace di alterare il clima planetario continuando a pensare con gli stessi strumenti con cui un tempo si cacciava. \`{E} esattamente la vergogna prometeica di Anders rovesciata: non siamo pi\`{u} inferiori alle nostre macchine, siamo inferiori alla portata delle nostre stesse azioni, incapaci di provare per i rischi che generiamo la paura proporzionata che Jonas ci chiedeva di coltivare deliberatamente, e incapaci di sentirci vincolati, come voleva Burke, da un patto con chi verr\`{a} dopo di noi e che non ha ancora voce per reclamarlo.

Siamo, in un certo senso, esseri capaci di un potere quasi assoluto sul pianeta ma ancora governati da un sistema di allarme progettato per la savana. Guardiamo il pianeta scaldarsi attraverso lo schermo di uno smartphone, e il nostro cervello pi\`{u} antico continua a sussurrarci che troveremo sempre una soluzione, che la tecnologia ci salver\`{a}, che il nostro gruppo ha ragione. Forse \`{e} cos\`{i}. O forse \`{e} proprio questa l'ultima, pi\`{u} pericolosa illusione: credere che i bias che ci hanno salvato nel passato non possano, questa volta, costarci il futuro.